% Generated draft: requires mathematical and source review.
\documentclass[11pt]{article}
\usepackage[margin=25mm]{geometry}
\usepackage{fontspec}
\setmainfont{Latin Modern Roman}
\usepackage{amsmath,amssymb,mathtools,unicode-math}
\setmathfont{Latin Modern Math}
\usepackage{enumitem,needspace,xurl,hyperref}
\setlength{\parindent}{0pt}
\setlength{\parskip}{5pt}
\setlength{\emergencystretch}{3em}
\newcommand{\N}{\mathbb{N}}
\newcommand{\Z}{\mathbb{Z}}
\newcommand{\Q}{\mathbb{Q}}
\newcommand{\R}{\mathbb{R}}
\newcommand{\C}{\mathbb{C}}
\newcommand{\F}{\mathbb{F}}
\newcommand{\A}{\mathbb{A}}
\newcommand{\PP}{\mathbb{P}}
\newcommand{\E}{\mathbb{E}}
\newcommand{\eps}{\varepsilon}
\newcommand{\dd}{\,\mathrm{d}}
\begin{document}
\section*{UnsolvedMath ID: 30003419}
\textbf{Source problem code:} catalogue entry
\subsection{Definitions and mathematical statement}
The supplied record describes a family of consistency questions referring to the report's Theorems 1--4. Those theorem statements, the affected cardinal invariants, and the parameter hypotheses are not included, so a complete quantified consistency assertion cannot be recovered from this record alone.

A cardinal \(\delta\) is supercompact if, for every cardinal \(\lambda\geq\delta\), there is an elementary embedding \(j:V\to M\) into a transitive class with critical point \(\delta\), with \(j(\delta)>\lambda\), and with \(M\) closed under \(\lambda\)-indexed sequences from \(V\). Here \(V\) denotes the ground model. For the source's unspecified parameter \(\kappa\), define, when the indicated set is nonempty,
\[
\theta^V=\min\{\delta>\kappa:V\models\text{``}\delta\text{ is supercompact''}\}.
\]
The record does not specify the domain or hypotheses for \(\kappa\).

Let \(\mathcal P\) stand schematically for an intended pattern of cardinal invariants from the report's Theorems 1--4, and let \((\mathfrak x_i)_{i\in I}\) stand for the invariants affected by the corresponding Boolean-ultrapower construction. Neither \(\mathcal P\) nor \(I\) is supplied; these symbols are placeholders, not reconstructed theorem statements.

The questions reported are whether the intended consistency patterns can be established without assuming a supercompact cardinal, or whether the affected invariants can consistently receive smaller, more accessible values than those produced by the reported constructions. The corrected formulation describes the latter direction using values below the first ground-model supercompact above \(\kappa\). If \(W\) denotes a resulting model and that ground-model cardinal exists, the indicated upper-bound condition would read
\[
\mathfrak x_i^W<\theta^V\qquad(i\in I),
\]
with the comparison made between ordinals. This display records only the stated bound; it does not specify the missing invariant pattern, the intended smaller values, or the required preservation properties. In particular, the record does not say which patterns or parameter choices are quantified over, or whether every affected invariant must satisfy the bound simultaneously.

The two directions must be distinguished: proving consistency without a supercompact-cardinal assumption does not itself mean requiring the resulting model to have no supercompact cardinals. Moreover, \(\theta^V\) is undefined if there is no ground-model supercompact above \(\kappa\). The original formulation asks generally for smaller values, whereas the corrected formulation additionally mentions weakening the large-cardinal assumption and the specific ground-model bound. Their exact relationship requires the missing theorem statements.
\subsection{Short English statement}
For the cardinal-invariant patterns in the report's Theorems 1--4, ask whether the supercompact-cardinal assumption can be avoided, or whether the invariants can consistently have smaller values, described in the corrected record as below the first ground-model supercompact above \(\kappa\). The theorems, invariants, parameter hypotheses, and intended smaller values are missing, so this remains a qualified question schema requiring mathematical review.
\subsection{Sources}
\begin{itemize}
\item Supplied local UnsolvedMath registry record, permanent ID \(30003419\), title: Reducing Cardinal Invariants from Boolean Ultrapower Constructions. Exact source problem code: \texttt{catalogue entry}. Collection: Oberwolfach Reports Open Problems. No individual proposer is supplied; the proposed year is recorded as 2017.
\item Supplied report citation: Set Theory (2018), Oberwolfach Reports, DOI 10.4171/owr/2017/11, pp. 527--589. \url{https://doi.org/10.4171/owr/2017/11}. The background describes Oberwolfach Reports as documenting mathematical workshops and their open-problem discussions. The record refers to Theorems 1--4 without supplying their text or a precise location for this question.
\item Supplied literature source, labelled OWR Set Theory report: \url{https://ems.press/content/serial-article-files/46674}.
\item Literature assessment dated exactly 2026-08-21: the supplied curation labels the question open, describes the extracted passage as a broad collection of Boolean-ultrapower questions, and says that later work gives related consistency constructions but no general determination of the requested smaller invariant values. This is supplied dated curation, not an independently verified assessment. No references identifying that later work are supplied. The background explicitly requests independent verification of time-sensitive status and cited claims.
\item The registry labels the statement \texttt{corrected\_verified} and reports that the official report was checked directly and that the clean statement restored definitions while preserving the target. These are supplied verification claims; no source visit or independent mathematical verification is asserted here.
\end{itemize}
\end{document}
